% =====================================================================
% Abstract
% File: paper/frontmatter/01-abstract.tex
% =====================================================================

\begin{abstract}
We develop a canonical framework for localized trace formulae on finite-area hyperbolic surfaces, with particular emphasis on stability under auxiliary analytic choices.  Starting from the Selberg trace formula, we construct a normalized local trace operator in which the identity sector is separated explicitly and all remaining non-canonical dependence is recorded in an additive error ledger.

The framework treats spectral localization, microlocal insertion, truncation, continuous-spectrum regularization, and refinement of partitions as controlled operations rather than informal modifications of a global trace identity.  For admissible test functions and localization data, the normalized trace decomposes into identity, hyperbolic, continuous-spectrum, and ledger terms.  We prove that admissible changes of cutoff profiles, quantization conventions, truncation schemes, and microlocal refinements affect the normalized trace only through explicit identity corrections and ledger-bounded remainders.

The resulting formalism provides a reusable local trace object suitable for quantitative spectral analysis.  In particular, we formulate power-saving regimes for short spectral windows, derive local Weyl-type consequences, and isolate trace-to-zeta interfaces in which determinant or explicit-formula constructions may be studied without hiding localization errors.  The emphasis throughout is on separating genuine spectral information from artifacts introduced by localization.

This work is intended as a foundational analytic layer: it does not identify the zeros of the Riemann zeta function with the spectrum of a hyperbolic surface, nor does it claim a critical-line theorem.  Instead, it supplies a stable local trace mechanism and an explicit ledger calculus that can be inserted into later determinant, zeta, or positivity problems when the corresponding external spectral object has been constructed independently.
\end{abstract}

\bigskip

\noindent\textbf{Keywords.} \PaperKeywords.

\medskip

\noindent\textbf{Mathematics Subject Classification (2020).} \PaperMSC.
